\section{Conclusion}
\label{sec:conclusion}

We studied whether recursive LLM ecosystems exhibit a minimum viable population threshold and how to define an individual agent operationally. Using a four-generation recursive inheritance proxy with six conditions, we found no clean threshold in which larger populations reliably prevented collapse. Instead, every condition showed a strong generation-1 truthfulness shock, followed by partial recovery. Our individuality assay gave a concrete answer to the population-definition question: agents count as distinct when between-agent distance exceeds within-agent stochastic variance. By that criterion, role-diverse prompts create distinct individuals, but that distinctness did not protect the ecosystem.

The main takeaway is straightforward. In this bounded setting, recursive stability depends more on inheritance quality and preserved human grounding than on nominal population count alone. Prompt-level role diversity is measurable, but it is not the form of diversity that prevents collapse.

Future work should replace the inheritance proxy with actual across-generation fine-tuning, introduce grounded diversity through retrieval or human refresh, expand the population design to include model-family diversity, and validate failure cases with stronger truthfulness judgments.
